\heading{高等数学A（II）-模拟题1-期末}

\noteinfo{题目和答案由AI生成。}

\begin{question}
求下列初值问题的解：
\begin{enumerate}[label=（\arabic*）,leftmargin=2.8em,itemsep=0.2em]
\item $y'-2y=e^{2x},\quad y(0)=1$；
\item $y'=y^{1/3},\quad y(0)=0$；
\item $y''-2y'+y=xe^x,\quad y(0)=0,\ y'(0)=1$。
\end{enumerate}
\begin{solution}

\begin{enumerate}[label=（\arabic*）,leftmargin=2.8em,itemsep=0.2em]
\item 用积分因子 $e^{-2x}$，
\[
(e^{-2x}y)'=1\Longrightarrow e^{-2x}y=x+C.
\]
由 $y(0)=1$ 得 $C=1$，故
\ansmath{y=(x+1)e^{2x}.}
\item 方程可分离：
\[
 y^{-1/3}dy=dx\Longrightarrow \frac32 y^{2/3}=x+C.
\]
因右端在 $y=0$ 处不满足 Lipschitz 条件，初值问题不唯一。典型解为
\[
 y\equiv0,
\]
以及对任意 $c\ge0$，
\[
 y(x)=\begin{cases}
 0,&x\le c,\\
 \bigl(\frac23(x-c)\bigr)^{3/2},&x\ge c.
 \end{cases}
\]
\item 设 $y=e^xv$，则
\[
 y''-2y'+y=e^xv''=xe^x,
\]
故 $v''=x$。积分得
\[
 v=\frac{x^3}{6}+C_1x+C_2.
\]
于是
\[
 y=e^x\left(\frac{x^3}{6}+C_1x+C_2\right).
\]
由初值 $y(0)=0$ 得 $C_2=0$；再算
\[
 y'(0)=C_1=1.
\]
故
\ansmath{y=e^x\left(x+\frac{x^3}{6}\right).}
\end{enumerate}
\end{solution}
\end{question}

\begin{question}
讨论函数项级数
\[
\sum_{n=1}^{\infty}\frac{\sin nx}{n^p}
\]
在下列区间上的一致收敛性：
\begin{enumerate}[label=（\arabic*）,leftmargin=2.8em,itemsep=0.2em]
\item $[\delta,2\pi-\delta]$，其中 $0<\delta<\pi$；
\item $[0,2\pi]$。
\end{enumerate}
\begin{solution}

\begin{enumerate}[label=（\arabic*）,leftmargin=2.8em,itemsep=0.2em]
\item 在 $[\delta,2\pi-\delta]$ 上，部分和
\[
\sum_{k=1}^{N}\sin kx=\frac{\sin\frac{Nx}{2}\sin\frac{(N+1)x}{2}}{\sin\frac{x}{2}}
\]
由于 $\sin(x/2)$ 远离 $0$，故部分和一致有界。再结合 $1/n^p\downarrow0$，由 Dirichlet 判别法知：
\ansmath{\text{对任意 }p>0,\ \sum\frac{\sin nx}{n^p}\text{ 在 }[\delta,2\pi-\delta]\text{上一致收敛。}}
\item 在 $[0,2\pi]$ 上，当 $p>1$ 时由 Weierstrass 判别法一致收敛；当 $0<p\le1$ 时不一致收敛，因为靠近 $x=0$ 时部分和不能一致有界，且点 $x=0$ 成为 Dirichlet 判别法失败的位置。故
\ansmath{[0,2\pi]\text{ 上一致收敛当且仅当 }p>1.}
\end{enumerate}
\end{solution}
\end{question}

\begin{question}
求函数 $f(x)=|x|$ 在 $[-\pi,\pi]$ 上的 Fourier 级数，并由此求和
\[
\sum_{n=1}^{\infty}\frac{1}{(2n-1)^2}.
\]

\begin{solution}

$f$ 为偶函数，仅含余弦项。先算
\[
a_0=\frac{2}{\pi}\int_0^{\pi}x\,dx=\pi.
\]
对 $n\ge1$，
\[
a_n=\frac{2}{\pi}\int_0^{\pi}x\cos nx\,dx
=\frac{2}{\pi}\frac{(-1)^n-1}{n^2}.
\]
故偶数项为零，奇数项为
\[
a_{2m-1}=-\frac{4}{\pi(2m-1)^2}.
\]
因此
\ansmath{|x|=\frac{\pi}{2}-\frac{4}{\pi}\sum_{m=1}^{\infty}\frac{\cos(2m-1)x}{(2m-1)^2},\qquad -\pi<x<\pi.}
取 $x=0$，得
\[
0=\frac{\pi}{2}-\frac{4}{\pi}\sum_{m=1}^{\infty}\frac{1}{(2m-1)^2},
\]
故
\ansmath{\sum_{m=1}^{\infty}\frac{1}{(2m-1)^2}=\frac{\pi^2}{8}.}
\end{solution}
\end{question}

\begin{question}
设
\[
 f(x)=\ln(1+x+x^2).
\]
求其在 $x=0$ 处的幂级数展开，并给出收敛半径。
\begin{solution}

注意
\[
1+x+x^2=\frac{1-x^3}{1-x},\qquad x\ne1.
\]
因此在 $|x|<1$ 内
\[
 f(x)=\ln(1-x^3)-\ln(1-x).
\]
用
\[
\ln(1-t)=-\sum_{n=1}^{\infty}\frac{t^n}{n},\qquad |t|<1,
\]
得
\[
 f(x)=\sum_{n=1}^{\infty}\frac{x^n}{n}-\sum_{n=1}^{\infty}\frac{x^{3n}}{n}.
\]
故展开为
\ansmath{f(x)=x+\frac{x^2}{2}-\frac{2x^3}{3}+\frac{x^4}{4}+\frac{x^5}{5}-\frac{x^6}{3}+\cdots}
更紧凑地写成
\ansmath{f(x)=\sum_{n=1}^{\infty}\frac{x^n}{n}-\sum_{n=1}^{\infty}\frac{x^{3n}}{n}.}
由于最近奇点来自 $1+x+x^2=0$ 的根 $e^{\pm 2\pi i/3}$，模长均为 $1$，故收敛半径
\ansmath{R=1.}
\end{solution}
\end{question}

\begin{question}
讨论级数
\[
\sum_{n=1}^{\infty}\frac{(-1)^n}{n^x+n\ln(n+1)}
\]
关于参数 $x\in\mathbb R$ 的收敛域与绝对收敛域。
\begin{solution}

对任意固定 $x\in\mathbb R$，分母均趋于 $+\infty$，且
\[
 a_n(x)=\frac{1}{n^x+n\ln(n+1)}
\]
最终单调下降到 $0$，故由 Leibniz 判别法，级数对一切 $x$ 收敛。

再看绝对收敛：
\[
\sum_{n=1}^{\infty}\frac{1}{n^x+n\ln(n+1)}.
\]
\begin{enumerate}[label=（\arabic*）,leftmargin=2.8em,itemsep=0.2em]
\item 若 $x>1$，分母主导项为 $n^x$，与 $\sum n^{-x}$ 同阶，绝对收敛；
\item 若 $x=1$，与 $\sum 1/(n\ln n)$ 同阶，发散；
\item 若 $x<1$，分母主导项为 $n\ln n$，仍与 $\sum 1/(n\ln n)$ 同阶，发散。
\end{enumerate}
故
\ansmath{\text{收敛域为 }\mathbb R,\qquad \text{绝对收敛域为 }(1,\infty).}
\end{solution}
\end{question}

\begin{question}
设参数积分
\[
F(a)=\int_0^{\infty}e^{-ax}\frac{\sin x}{x}\,dx,\qquad a>0.
\]
求 $F(a)$。
\begin{solution}

对 $a>0$ 可对参数求导：
\[
F'(a)=-\int_0^{\infty}e^{-ax}\sin x\,dx=-\frac{1}{1+a^2}.
\]
故
\[
F(a)=-\arctan a+C.
\]
又当 $a\to\infty$ 时，由支配收敛定理 $F(a)\to0$，于是
\[
0=-\frac\pi2+C\Longrightarrow C=\frac\pi2.
\]
故
\ansmath{F(a)=\frac\pi2-\arctan a=\arctan\frac1a.}
\end{solution}
\end{question}


